\subsection{Edge Bundle Construction with the Robot Dynamics Model}\label{sec:methods-eb_construction}

The backbone of our planner lies in constructing a discrete approximation of the robot's state and control spaces, as a result capturing the nuances of real-world dynamics including nonlinearities like joint friction. Defined as a bundle of edges $\mathcal{E}$, this approximation is a collection of valid, disconnected, and randomly generated forward rollouts of the robot dynamics model over a variable duration $\Delta t$. This method of edge generation with a sampled duration is known as Monte Carlo propagation and ensures the preservation of asymptotic optimality in kinodynamic planning algorithms \cite{kunz2015kinodynamic}.
In order to enable uniform and rapid exploration of the control space, we propagate a sequence of randomly-sampled joint velocity actions, $\boldsymbol{\dot{q}} = (\dot{q}_0, \cdots, \dot{q}_f) \in \mathcal{U}$, from a randomly sampled state, $q \in \mathcal{Q}$. In addition to providing a broader local reachability range (and as a result expanding the explored volume in the state space) from a given state, varying joint velocities over a single propagation leads to longer-duration edges as it prevents the robot from becoming locally ``locked'' along fixed controls.
Moreover, it is necessary that these Monte Carlo propagations are collision-free and obey the Euler-Lagrange robot dynamics model $\mathcal{D}$ to ensure validity in real-world execution. Our approach to bundle generation can be extended to other robot embodiments as long as a valid dynamics model is provided.

As an example, we focus on one particular embodiment and use an analytically-derived and empirically-tuned dynamics model for the 7 degrees-of-freedom (DoF) Franka Emika Panda arm \cite{Gaz2019}. This analytical model is defined as
$\tau(q, \dot{q}, \ddot{q}) = M(q)\ddot{q} + C(q, \dot{q})\dot{q} + g(q) + f(\dot{q})$,
where $q$, $\dot{q}$, $\ddot{q}$, and $\tau$ are $7\times1$ vectors representing the robot's joint angles, velocities, accelerations, and torques, respectively. $M$ is a symmetric, positive-definite $7\times7$ joint inertia matrix, $C$ is a $7\times1$ matrix that captures the Coriolis and centrifugal effects caused by robot motion, $g$ is a $7\times1$ vector encapsulating the torque required by each motor to counteract gravitational forces, and $f$ is a $7\times1$ vector that quantifies the torques necessary to counteract friction at the joints. The model is used to verify whether the propagated joint velocity sequence $\boldsymbol{\dot{q}}$ obeys the robot torque limits.

Furthermore, validating an edge involves three classes of checks: collision checks, robot state limit checks, and continuity checks (\cref{algo:lazyboe}, Lines 20-22).
Acceptable ranges for the state and control variables and additionally, dynamics quantities are defined as $q_{min} < q < q_{max}$, $\dot{q}_{min} < \dot{q} < \dot{q}_{max}$, $\ddot{q}_{min} < \ddot{q} < \ddot{q}_{max}$, and $\tau_{min} < \tau < \tau_{max}$.
In addition, the dynamics simulations must obey continuity constraints due to the stringent control requirements of this robot.
For a given static environment, edges are validated to be free of self-collision and environmental collisions via geometric collision-checking algorithms. This incurs a one-time cost since the edges can be used lazily for a significant portion of our motion planning process.
We next explain how we can compute useful quantities from this generated edge bundle to enable lazy kinodynamic planning.
